\chapter{Conception de la base de donn\'ees}

Ce chapitre pr\'esente la conception du mod\`ele de donn\'ees du syst\`eme Andon, depuis l'analyse des besoins en donn\'ees jusqu'\`a l'impl\'ementation SQL dans PostgreSQL.

\section{M\'ethodologie de conception}

\subsection{Approche adopt\'ee}

La conception de la base de donn\'ees suit une approche descendante :

\begin{enumerate}
    \item Identification des entit\'es et de leurs attributs.
    \item D\'efinition des relations entre les entit\'es.
    \item Normalisation du mod\`ele (3FN minimum).
    \item Impl\'ementation du sch\'ema SQL.
    \item Ajout des index et optimisations.
\end{enumerate}

\subsection{Entit\'es principales}

L'analyse des besoins a permis d'identifier les entit\'es suivantes :

\begin{itemize}
    \item \textbf{Ligne :} repr\'esente une ligne de production de l'atelier.
    \item \textbf{Evenement :} repr\'esente un \'ev\'enement Andon (panne, anomalie).
    \item \textbf{Utilisateur :} repr\'esente un utilisateur du syst\`eme.
    \item \textbf{Notification :} repr\'esente une notification envoy\'ee.
    \item \textbf{Capteur :} repr\'esente un n\oe{}ud ESP32 connect\'e.
\end{itemize}

\section{Mod\`ele Entit\'e-Association}

\subsection{Diagramme EA}

Le mod\`ele Entit\'e-Association d\'ecrit les entit\'es, leurs attributs et les relations entre elles :

\begin{itemize}
    \item \textbf{Ligne -- Evenement} : une ligne peut g\'en\'erer plusieurs \'ev\'enements (relation 1:N).
    \item \textbf{Utilisateur -- Evenement} : un utilisateur peut cr\'eer ou g\'erer plusieurs \'ev\'enements (relation 1:N).
    \item \textbf{Ligne -- Capteur} : une ligne est associ\'ee \`a un capteur (relation 1:1).
    \item \textbf{Evenement -- Notification} : un \'ev\'enement peut d\'eclencher plusieurs notifications (relation 1:N).
\end{itemize}

\figplaceholder{Diagramme Entit\'e-Association du syst\`eme Andon}{fig:diagramme_ea}

\subsection{D\'efinition des attributs}

\subsubsection{Entit\'e Ligne}

\begin{table}[htbp]
\centering
\renewcommand{\arraystretch}{1.3}
\begin{tabular}{l l l}
\toprule
\textbf{Attribut} & \textbf{Type} & \textbf{Description} \\
\midrule
id & SERIAL (PK) & Identifiant unique \\
nom & VARCHAR(100) & Nom de la ligne \\
numero & INTEGER & Num\'ero de la ligne \\
statut & VARCHAR(20) & Statut actuel \\
date\_creation & TIMESTAMP & Date de cr\'eation \\
\bottomrule
\end{tabular}
\caption{Attributs de l'entit\'e Ligne}
\label{tab:attributs_ligne}
\end{table}

\subsubsection{Entit\'e Evenement}

\begin{table}[htbp]
\centering
\renewcommand{\arraystretch}{1.3}
\begin{tabular}{l l l}
\toprule
\textbf{Attribut} & \textbf{Type} & \textbf{Description} \\
\midrule
id & SERIAL (PK) & Identifiant unique \\
ligne\_id & INTEGER (FK) & R\'ef\'erence \`a la ligne \\
type & VARCHAR(50) & Type d'\'ev\'enement \\
statut & VARCHAR(20) & En cours / R\'e-solu \\
description & TEXT & Description d\'etail\'ee \\
date\_debut & TIMESTAMP & Heure de d\'etection \\
date\_fin & TIMESTAMP & Heure de r\'e-solution \\
technicien\_id & INTEGER (FK) & Technicien assign\'e \\
compte\_rendu & TEXT & Compte-rendu d'intervention \\
\bottomrule
\end{tabular}
\caption{Attributs de l'entit\'e Evenement}
\label{tab:attributs_evenement}
\end{table}

\section{Mod\`ele relationnel}

\subsection{Sch\'ema relationnel}

Le passage du mod\`ele EA au mod\`ele relationnel produit les tables suivantes :

\begin{lstlisting}[caption={Sch\'ema relationnel},label={lst:schema_relational}]
Ligne (id, nom, numero, statut, date_creation)
Evenement (id, ligne_id, type, statut, 
           description, date_debut, date_fin, 
           technicien_id, compte_rendu)
Utilisateur (id, nom, prenom, email, role, 
             mot_de_passe, date_creation)
Notification (id, evenement_id, destinataire_id, 
              type, contenu, date_envoi, statut)
Capteur (id, ligne_id, mac_address, 
         firmware_version, date_installation)
\end{lstlisting}

\subsection{R\`egles d'int\'egrit\'e}

Les contraintes d'int\'egrit\'e suivantes sont appliqu\'ees :

\begin{itemize}
    \item \textbf{Cl\'es primaires :} chaque table poss\`ede un identifiant unique auto-incr\'ement\'e (SERIAL).
    \item \textbf{Cl\'es \'etrang\`eres :} les r\'ef\'erences entre tables sont contr\^ol\'ees par des contraintes FOREIGN KEY.
    \item \textbf{NOT NULL :} les attributs critiques ne peuvent \^etre nuls.
    \item \textbf{UNIQUE :} l'adresse e-mail des utilisateurs doit \^etre unique.
    \item \textbf{CHECK :} les statuts doivent respecter les valeurs d\'efinies.
\end{itemize}

\section{Impl\'ementation SQL}

\subsection{Cr\'eation des tables}

\begin{lstlisting}[language=SQL,caption={Cr\'eation des tables principales},label={lst:create_tables}]
CREATE TABLE lignes (
    id SERIAL PRIMARY KEY,
    nom VARCHAR(100) NOT NULL,
    numero INTEGER UNIQUE NOT NULL,
    statut VARCHAR(20) DEFAULT 'active' 
        CHECK (statut IN (
            'active','inactive','en_panne'
        )),
    date_creation TIMESTAMP 
        DEFAULT CURRENT_TIMESTAMP
);

CREATE TABLE utilisateurs (
    id SERIAL PRIMARY KEY,
    nom VARCHAR(100) NOT NULL,
    prenom VARCHAR(100) NOT NULL,
    email VARCHAR(255) UNIQUE NOT NULL,
    role VARCHAR(50) NOT NULL
        CHECK (role IN (
            'admin','technicien','chef','operateur'
        )),
    mot_de_passe VARCHAR(255) NOT NULL,
    date_creation TIMESTAMP 
        DEFAULT CURRENT_TIMESTAMP
);

CREATE TABLE evenements (
    id SERIAL PRIMARY KEY,
    ligne_id INTEGER NOT NULL 
        REFERENCES lignes(id),
    type VARCHAR(50) NOT NULL,
    statut VARCHAR(20) DEFAULT 'en_cours'
        CHECK (statut IN ('en_cours','resolu')),
    description TEXT,
    date_debut TIMESTAMP 
        DEFAULT CURRENT_TIMESTAMP,
    date_fin TIMESTAMP,
    technicien_id INTEGER 
        REFERENCES utilisateurs(id),
    compte_rendu TEXT
);

CREATE TABLE notifications (
    id SERIAL PRIMARY KEY,
    evenement_id INTEGER 
        REFERENCES evenements(id),
    destinataire_id INTEGER 
        REFERENCES utilisateurs(id),
    type VARCHAR(50) NOT NULL,
    contenu TEXT,
    date_envoi TIMESTAMP 
        DEFAULT CURRENT_TIMESTAMP,
    statut VARCHAR(20) DEFAULT 'envoyee'
);
\end{lstlisting}

\subsection{Index de performance}

Pour garantir des performances optimales, les index suivants sont cr\'e\'es :

\begin{lstlisting}[language=SQL,caption={Index de performance},label={lst:indexes}]
CREATE INDEX idx_evenements_ligne 
    ON evenements(ligne_id);
CREATE INDEX idx_evenements_date 
    ON evenements(date_debut);
CREATE INDEX idx_evenements_statut 
    ON evenements(statut);
CREATE INDEX idx_evenements_date_ligne 
    ON evenements(date_debut, ligne_id);
CREATE INDEX idx_notifications_dest 
    ON notifications(destinataire_id);
\end{lstlisting}

\subsection{Donn\'ees de test}

Des donn\'ees de test ont \'et\'e ins\'er\'ees pour valider le syst\`eme :

\begin{lstlisting}[language=SQL,caption={Insertion des lignes de test},label={lst:insert}]
INSERT INTO lignes (nom, numero, statut) VALUES
('Ligne T1', 1, 'active'),
('Ligne T2', 2, 'active'),
('Ligne T3', 3, 'active'),
('Ligne T4', 4, 'active'),
('Ligne T5', 5, 'active'),
('Ligne T6', 6, 'active'),
('Ligne T7', 7, 'active'),
('Ligne T8', 8, 'active');
\end{lstlisting}

\section{Gestion des volumes de donn\'ees}

\subsection{Strat\'egie de r\'etention}

\'Etant donn\'e le volume croissant d'\'ev\'enements g\'en\'er\'es, une strat\'egie de r\'etention est mise en place :

\begin{itemize}
    \item Les \'ev\'enements d\'etaill\'es sont conserv\'es pendant 12 mois.
    \item Les donn\'ees agr\'eg\'ees sont conserv\'ees ind\'efiniment.
    \item Un job de nettoyage p\'eriodique supprime les donn\'ees anciennes.
\end{itemize}

\subsection{Partitionnement temporel}

Pour les volumes importants, le partitionnement temporel est envisag\'e :

\begin{lstlisting}[language=SQL,caption={Partitionnement temporel},label={lst:partition}]
CREATE TABLE evenements (
    ...
) PARTITION BY RANGE (date_debut);

CREATE TABLE evenements_2026_q1 
    PARTITION OF evenements
    FOR VALUES FROM ('2026-01-01') 
    TO ('2026-04-01');
\end{lstlisting}

\section{Synth\`ese du chapitre}

La conception de la base de donn\'ees suit une approche rigoureuse allant du mod\`ele EA au sch\'ema SQL. Le mod\`ele comprend cinq entit\'es principales : Ligne, Evenement, Utilisateur, Notification et Capteur. L'impl\'ementation dans PostgreSQL int\`egre les contraintes d'int\'egrit\'e, les index de performance et une strat\'egie de gestion des volumes de donn\'ees adapt\'ee aux besoins du syst\`eme Andon.
